\section{Materials and methods}
\label{sec:method}

\Cref{fig:pipeline} summarises the pipeline. A frozen backbone produces a
per-epoch distribution; an HMM fitted to the training labels smooths it over
the recording; a nonconformity score turns the smoothed posterior into a
scalar level; and one of three calibration rules --- pooled split CP, a
subject-adaptive warm-up, or a per-patient adaptive stream --- turns that
level into a prediction set. We describe each step and state the assumption
it needs.

\begin{figure}[t]
  \centering
  \resizebox{\linewidth}{!}{% 
  \begin{tikzpicture}[
    box/.style={draw, rounded corners=2pt, align=center,
                minimum height=10mm, minimum width=20mm,
                inner sep=4pt, font=\small},
    cal/.style={box, fill=black!5, text width=31mm},
    arr/.style={-{Latex[length=1.6mm]}, thick},
    lab/.style={font=\scriptsize\itshape, text=black!70},
    note/.style={font=\scriptsize, align=center, text=black!75}]

    % First pass: one EEG epoch becomes a temporally smoothed posterior.
    \node[box] (x) at (0,0) {EEG epoch\\$X_t\in\bR^{3000}$};
    \node[box] (f) at (2.7,0) {frozen\\backbone $f$};
    \node[box] (p) at (5.4,0) {softmax\\$p_t$};
    \node[box] (h) at (8.5,0) {HMM\\forward--backward};
    \node[box] (g) at (11.8,0) {posterior\\$\gamma_t$};

    % Second pass: posterior, calibrated threshold, and prediction set.
    \node[box] (s) at (11.8,-1.8) {score\\$s_t(y)$};
    \node[box] (q) at (8.5,-1.8) {threshold\\$\wh q(\alpha_t)$};
    \node[box] (c) at (5.4,-1.8) {set\\$\cC_t$};

    % The two auxiliary inputs sit on their own rows and cannot overlap.
    \node[cal] (cal) at (7.6,-3.75)
      {calibration subjects\\$\{S_{j,t}\}$\\$\rightarrow$ pooled quantile};
    \node[cal] (fb) at (12.0,-3.75)
      {review: label\\$Y_t$ after $\cC_t$};

    \draw[arr] (x) -- (f);
    \draw[arr] (f) -- (p);
    \draw[arr] (p) -- (h);
    \draw[arr] (h) -- (g);
    \draw[arr] (g) -- (s);
    \draw[arr] (s) -- (q);
    \draw[arr] (q) -- (c);
    \draw[arr, dashed] (cal.north) -- (q.south);
    \draw[arr, dashed] (fb.north) -- (s.south);

    % Feedback takes a separate lower route, leaving the calibration box clear.
    \draw[arr, dashed] (c.south) -- (5.4,-5.1) -- (11.8,-5.1) -- (fb.south);
    \node[note] at (8.6,-5.45) {reviewed labels feed the adaptive stream};

    \node[lab, above=2mm of h] {transition matrix from training labels};
  \end{tikzpicture}}%
  \caption{Overview of the set-valued sleep-staging pipeline. A frozen
  backbone produces epoch probabilities; forward--backward inference uses
  transition probabilities estimated from training labels to form a smoothed
  posterior. The calibration subjects determine the threshold, and the
  reviewed stage label can update the per-patient adaptive threshold. Dashed
  paths are used only in the adaptive regime.}
  \label{fig:pipeline}
\end{figure}

\subsection{Data}
\label{sec:data}

\paragraph{Sleep-EDF}
We use the Sleep-EDF Expanded corpus from PhysioNet
\citep{kemp2000,goldberger2000}, single-channel Fpz--Cz EEG at
\SI{100}{\hertz}. Two cohorts are used. \emph{Sleep-EDF-20} is the
20-subject cohort (SC400--SC419) commonly used in the sleep-staging
literature: 39 recordings, \num{42368} epochs after cropping, of which
\num{42307} carry a stage label (the remaining epochs sit inside the cropped
window but fall between scored annotations). The
\emph{extended cohort} adds the remaining cassette recordings (SC400--SC482),
giving 78 subjects and 153 recordings, and is used for the
statistical-power analysis in \Cref{sec:results-cohort}.

\paragraph{Telemetry cohort}
Sleep-EDF also contains a telemetry subset (ST700--ST721, 22 subjects,
44 recordings) recorded with different hardware on a different patient
group. It is never used for training or calibration in the cross-cohort
experiment, so it provides a genuine hardware-and-population shift rather
than a resampling of the same distribution.

\paragraph{Preprocessing}
Each recording is band-pass filtered to \SIrange{0.5}{30}{\hertz}, segmented
into \SI{30}{\second} epochs, and $z$-scored per recording. A label is
assigned to an epoch when at least 50\% of it overlaps a scored annotation.
R\&K stage~4 is merged into N3 following the AASM rule
\citep{berry2012manual}, because keeping it as a separate class would create a
sixth category with a handful of epochs and no modern counterpart.

Each recording is cropped to the scored sleep period plus 30 minutes of wake
on either side. Keeping every annotated epoch instead yields a cohort of over
\num{112000} epochs that is $65\%$ wake, which matches no published
result and distorts the class priors; the crop reproduces the canonical
Sleep-EDF-20 cohort with priors W $19.6\%$, N1 $6.6\%$, N2
$42.1\%$, N3 $13.5\%$ and REM $18.2\%$. The crop is applied
identically to every cohort and is part of the released code. Subject-wise
folds, training, and the evaluation metrics are described in
\Cref{sec:setup}.

\subsection{Notation and problem setting}
\label{sec:notation}

Let $i$ index subjects, $r$ recordings and $t$ epochs. A subject is one
patient (one or two recordings), and we use the two words interchangeably.
Each epoch has an input
$X_t\in\bR^{3000}$ (30\,s of a single EEG channel at \SI{100}{\hertz}) and a
label $Y_t\in\cY=\{W,N1,N2,N3,REM\}$ with $|\cY|=5$. A backbone $f$, trained
on subjects disjoint from the test subject, outputs
$p_t=\operatorname{softmax}(f(X_t))\in\Delta^{4}$, where $\Delta^4$ is the
probability simplex. We are given labelled calibration data from subjects
disjoint from the test subject, and we must output a set
$\cC_t\subseteq\cY$ such that
\begin{equation}
\bP\big(Y_t\in\cC_t\big)\ \ge\ 1-\alpha
\label{eq:coverage}
\end{equation}
at a nominal level $\alpha$, while keeping $|\cC_t|$ small, because every
epoch with $|\cC_t|>1$ is an epoch a human has to look at.

Two remarks fix what \Cref{eq:coverage} does and does not say. First, the
probability is over the draw of calibration and test data, not over
subjects: it is a \emph{marginal} statement, and a predictor can satisfy it
while failing on a subset of patients. Second, in sleep staging the
natural exchangeable unit is arguably the patient rather than the epoch,
and the two differ exactly in the regime of interest, where patients are
heterogeneous.

We report, alongside coverage, the \emph{set size} $|\cC_t|$, the
\emph{singleton rate} $\bP(|\cC_t|=1)$, the \emph{review workload}
$\bP(|\cC_t|>1)$, the fraction of epochs with more than one admissible stage,
the \emph{empty rate} $\bP(|\cC_t|=0)$, which also returns an epoch to the
reviewer and which the workload does not count, and a \emph{temporal
consistency} statistic, the mean Jaccard
overlap of the sets at adjacent epochs,
\begin{equation}
\mathrm{TC}=\frac{1}{|\mathcal{P}|}\sum_{(t,t+1)\in\mathcal{P}}
\frac{|\cC_t\cap\cC_{t+1}|}{|\cC_t\cup\cC_{t+1}|},
\end{equation}
where $\mathcal{P}$ contains adjacent pairs inside a recording, so that
recording boundaries are never counted. $\mathrm{TC}=1$ when neighbouring
sets are identical and $0$ when they are disjoint.

\subsection{Split conformal prediction}
\label{sec:split-cp}

Given a nonconformity score $s(x,y)$ for which smaller values indicate better
agreement between input $x$ and candidate label $y$, split CP computes
$S_j=s(X_j,Y_j)$ on $n$ calibration epochs and takes the
$\lceil(n+1)(1-\alpha)\rceil/n$ empirical quantile,
\begin{equation}
\wh q=\operatorname{Quantile}_{\lceil(n+1)(1-\alpha)\rceil/n}\big(S_1,\dots,S_n\big),
\qquad
\cC(x)=\{y\in\cY:\ s(x,y)\le\wh q\}.
\label{eq:split-cp}
\end{equation}
If the calibration and test epochs are exchangeable, then
$\bP(Y\in\cC(X))\ge 1-\alpha$ and $\bP(Y\in\cC(X))\le 1-\alpha+1/(n+1)$: the
guarantee is finite-sample and distribution-free, and both bounds are tight
for the class of all distributions \citep{vovk2005,lei2018}.

\subsection{Nonconformity scores}
\label{sec:scores}

The \emph{adaptive prediction set} score \citep{romano2020} is
\begin{equation}
s_{\aps}(x,y)=\sum_{y'\in\cY}p_{y'}(x)\,\ones\{p_{y'}(x)\ge p_y(x)\},
\label{eq:aps}
\end{equation}
the cumulative probability mass of all classes at least as likely as $y$, so
that $s_{\aps}\in[0,1]$ and smaller is better. Its set is the smallest
collection of most-likely classes whose total mass reaches $\wh q$. RAPS
\citep{angelopoulos2021raps} adds a rank penalty $\lambda_{\mathrm{reg}}\max(0,\mathrm{rank}(y)-k_{\mathrm{reg}})$
to discourage a long tail of unlikely labels.

APS is heavy-tailed near one whenever the model is badly wrong, which forces
the calibration quantile up and inflates every set. For the temporal
component we therefore use the \emph{posterior deficit}
\begin{equation}
s_{\mathrm{def}}(x,y)=\max_{k\in\cY}\gamma_x(k)-\gamma_x(y),
\label{eq:deficit}
\end{equation}
where $\gamma_x$ is a smoothed posterior (\Cref{sec:method-smoothing}). This
score is bounded in $[0,1]$, is zero exactly at the maximum-posterior stage,
and measures how far a candidate is from the posterior's mode.

A \emph{class-conditional} (Mondrian) construction \citep{vovk2012} replaces
the pooled quantile by one quantile per stage,
\begin{equation}
\wh q_y=\operatorname{Quantile}_{\lceil(n_y+1)(1-\alpha)\rceil/n_y}
\big(\{S_j:\ Y_j=y\}\big),
\label{eq:mondrian}
\end{equation}
where $n_y$ is the number of calibration epochs of stage $y$, and sets
$\cC(x)=\{y\in\cY: s(x,y)\le\wh q_y\}$. The guarantee then holds within each
stage, $\bP(Y\in\cC(X)\mid Y=y)\ge 1-\alpha$, at the price of a smaller
calibration set per class; it is the baseline we use whenever a comparison
needs stage-level coverage rather than set size.

\subsection{Adaptive conformal inference}
\label{sec:aci}

When exchangeability fails, adaptive conformal inference \citep{gibbs2021}
updates the target miscoverage after each observed label,
\begin{equation}
\alpha_{t+1}=\alpha_t+\gamma_t\big(\alpha-\ones\{Y_t\notin\cC_t\}\big),
\qquad
\cC_t=\{y:\ s(X_t,y)\le \wh q(\alpha_t)\},
\label{eq:aci}
\end{equation}
where $\wh q(\alpha_t)$ is the calibration quantile at level $1-\alpha_t$.
Along an arbitrary stream, a bounded step keeps the long-run miscoverage
within a term governed by the step size rather than letting it drift with the
shift; \citet{angelopoulos2024} show that a decaying schedule with
$\sum_t\gamma_t=\infty$ and $\sum_t\gamma_t^2<\infty$ makes that term vanish,
which is the strongest form of the guarantee and the one the corollary below
assumes.
